\section{Benchmarks and results: Ingestion}
\label{sec:benchmarks-ingestion}

Ingestion was measured with JMH on the host: the work of a worker on one book, without the network of Project Gutenberg,
and apart from it the cost of coordinating that book through Hazelcast. The scripts, the instructions to repeat the
measure and the raw results are in \code{benchmarks/ingestion/}.

\subsection{Microbenchmarks with JMH}

The benchmarks live in the test code of \code{ingestion-service}, package \code{...ingestion.benchmark}, and run from the
benchmarks jar on the reference machine (Section~\ref{sec:consolidation}), with the same settings as Indexing and Search:
2 forks with a 4~GB heap, 5 warm-up and 10 measured iterations of 1~s each (20 samples), average time per operation with
its 99.9~\% error. The books are the 200 synthetic books of the dataset, the same bytes as the Indexing benchmarks, read
from the page cache. The benchmarks with threads take a round of 100 books, shared by 1, 2, 4 or 8 threads as the workers
of a replica share \code{ingest-queue}; every round writes into an empty datalake, and after it, unmeasured, the benchmark
checks that the 100 books are complete (or \code{DOWNLOADED} and in \code{index-queue}), so a wrong round fails instead of
giving a number.

\begin{table}[H]
\centering\small
\begin{tabularx}{\linewidth}{@{}Xrl@{}}
\toprule
\textbf{Benchmark} & \textbf{Time ($\mu$s)} & \textbf{Derived} \\
\midrule
Split a short book (8.0~KB raw) & 11.5 $\pm$ 1.0 & 695 MB/s \\
Split the median book (51.7~KB) & 79.3 $\pm$ 5.8 & 651 MB/s \\
Split a long book (287.8~KB) & 454 $\pm$ 15 & 634 MB/s \\
Save a book in the range datalake, per book, 1 thread & 235 $\pm$ 5 & 551 MB/s \\
\quad 2 threads & 142 $\pm$ 4 & $\times$1.66 \\
\quad 4 threads & 112 $\pm$ 4 & $\times$2.09 \\
\quad 8 threads & 121 $\pm$ 2 & $\times$1.93 \\
Ingest a book from the local source, per book, 1 worker & 473 $\pm$ 8 & 2,113 books/s \\
\quad 2 workers & 270 $\pm$ 4 & $\times$1.75 \\
\quad 4 workers & 175 $\pm$ 4 & $\times$2.70 \\
\quad 8 workers & 155 $\pm$ 2 & $\times$3.06 \\
Claim and complete a book in Hazelcast, per book, 1 worker & 1,631 $\pm$ 108 & 613 books/s \\
\quad 2 workers & 1,052 $\pm$ 66 & $\times$1.55 \\
\quad 4 workers & 581 $\pm$ 19 & $\times$2.81 \\
\quad 8 workers & 364 $\pm$ 15 & $\times$4.48 \\
\bottomrule
\end{tabularx}
\caption{JMH of Ingestion. \emph{Ingest} is \code{IngestionService.ingest}: read the raw file, split it, check that the
book is not in the datalake and save it (header, body and \code{\_READY}, each through a temporary file and an atomic
move). \emph{Claim and complete} are T2 and T3 of the contract, two transactions each, on a cluster of two members
started in the same JVM, the threads sharing one client as the workers of a replica do. The factors are the speed-up over
1 thread; the books of a round average 129~KB written.}
\label{tab:jmh-ingestion}
\end{table}

\textbf{Where the time of a book goes.} With one worker, ingesting a book from the local source takes 473~$\mu$s:
about 200~$\mu$s to split it (the splitter runs at 630 to 700~MB/s, linear in the size of the book), 235~$\mu$s to write
its three files, and the remaining 35~$\mu$s to read it from the page cache and check that it is not stored yet. Nothing is
forced to disk (the datalake does not \code{fsync}, Section~\ref{sec:readiness}), so the write measures the file system up
to the page cache.

\textbf{Threads.} The ingestion speeds up 2.7 times with 4 workers and 3.1 times with 8, on 6 physical cores. The write
stops at 4 threads and is slower with 8: the 100 books of a round, with consecutive ids, go to the same range folder, as
they do in the service, and every atomic move changes that one directory. The probable cause is that the file system
updates one directory one change at a time; it was not measured apart.
The splitting, pure CPU, keeps scaling, which is why the whole ingestion still gains a little at 8.

\textbf{Coordination.} Claiming and completing a book costs 1.6~ms with one worker: 3.4 times the ingestion of the book
itself. Each one is a transaction with a two-phase commit, so the time is round trips rather than CPU, and concurrent
workers overlap them: with 8 workers the cost per book falls to 0.36~ms ($\times$4.5), better than any other part.
With the default of 4 workers, coordination takes 581~$\mu$s per book against 175~$\mu$s of work: three quarters of the
time of a book with the local source. In this benchmark the members are in the same JVM and talk over loopback, so it is a
lower bound of the network cost and an upper bound of the CPU they compete for; the queue is never empty, so \code{claim}
contains no wait for an id.

\textbf{What it means for the pipeline.} With the local source, one replica with 4 workers is bounded by Hazelcast at
about 1,700 books per second in these conditions. That is two orders of magnitude above Indexing, which does 15 books/s on
2,000 books (Section~\ref{sec:benchmarks-indexing}), so Ingestion is never the bottleneck of the pipeline. With Project
Gutenberg as the source the download is much longer than both parts and the throttling of the site sets the pace; that is
why the measures use the local source.

\textbf{Not measured.} The integration experiments of Ingestion with the real services (1, 2 and 4 replicas, the 10,000
books and the coordination cost inside Docker) were prepared (\code{run-integration.sh}, in the branch
\code{test/CF-13-integration}) but not run: the group left them as not measured. Without them there is no
measure of the contention of several replicas on the single partition of each queue
(Section~\ref{sec:decisions-coordination}), nor of the overhead of Docker on the coordination; the services expose the
needed metrics in \code{GET /status} (\code{coordination\_cost}, CF-12) for whoever runs them.
